\documentclass[11pt, a4paper]{article}
\usepackage[a4paper, margin=2.5cm, headheight=14pt]{geometry}
\usepackage[utf8]{inputenc}
\usepackage[T1]{fontenc}
\usepackage[spanish, es-tabla]{babel}

\usepackage{amsmath}
\usepackage{amssymb}
\usepackage{booktabs}
\usepackage{graphicx}
\usepackage{xcolor}
\usepackage{enumitem}
\usepackage{fancyhdr}
\usepackage{titlesec}
\usepackage[colorlinks=true, linkcolor=primary, urlcolor=primary, citecolor=primary]{hyperref}

% Definición de paleta de colores institucional / académica
\definecolor{primary}{RGB}{15, 76, 129}      % Azul clásico institucional
\definecolor{secondary}{RGB}{90, 105, 120}   % Gris azulado para subtítulos/líneas
\definecolor{darktext}{RGB}{33, 37, 41}      % Texto oscuro principal
\definecolor{lightbg}{RGB}{248, 249, 250}    % Fondo claro para cajas/resumen

% Formato estilizado de Secciones
\titleformat{\section}
  {\color{primary}\normalfont\Large\bfseries}
  {\thesection.}{0.8em}{}
  [\vspace{0.2em}\color{secondary}\hrule height 0.8pt]

\titleformat{\subsection}
  {\color{secondary}\normalfont\large\bfseries}
  {\thesubsection}{0.6em}{}

% Encabezados y pies de página profesionales
\pagestyle{fancy}
\fancyhf{}
\rhead{\small\color{secondary}\itshape Tecnología Electrónica -- Amplificador TDA2030}
\lhead{\small\color{secondary}\leftmark}
\rfoot{\small\color{secondary}Página \thepage}
\lfoot{\small\color{secondary}Informe de Laboratorio}
\renewcommand{\headrulewidth}{0.4pt}
\renewcommand{\footrulewidth}{0.4pt}

\begin{document}

% ==============================================================================
% PORTADA (Omitida aquí por brevedad, usa la tuya anterior)
% ==============================================================================

\newpage
\pagenumbering{roman}

\begin{abstract}
\noindent Este documento constituye el informe técnico final correspondiente al desarrollo del amplificador de audio de potencia basado en el circuito integrado monolítico \textbf{TDA2030}. Se presenta el análisis teórico en Clase AB, la caracterización de la red de realimentación, el flujo completo de diseño CAD/EDA utilizando la herramienta KiCad, los resultados del análisis transitorio y de respuesta en frecuencia mediante simulación SPICE, y finalmente el procedimiento de montaje físico, pruebas de laboratorio y verificación de desempeño.
\end{abstract}

\vspace{1.5em}

% ==============================================================================
% ÍNDICE SEGURO (Sin tocloft)
% ==============================================================================
{
  \hypersetup{linkcolor=darktext} % Opcional: para que los enlaces del índice sean sobrios
  \renewcommand{\contentsname}{\large\bfseries\color{primary}ÍNDICE GENERAL}
  \tableofcontents
}

\newpage
\pagenumbering{arabic}

% ==============================================================================
% CONTENIDO DEL INFORME
% ==============================================================================
\section{Introducción y Objetivos}
El propósito de este trabajo práctico es aplicar las metodologías de diseño y fabricación de circuitos electrónicos mediante el desarrollo de un amplificador de audio operacional de potencia. Los objetivos principales incluyen:
\begin{itemize}[label=\small\color{primary}$\blacksquare$]
    \item Comprender las etapas fundamentales de amplificación de señal en baja frecuencia.
    \item Diseñar la placa de circuito impreso (PCB) en KiCad a partir del esquema eléctrico del datasheet.
    \item Realizar la simulación circuital para predecir el comportamiento en frecuencia y potencia.
    \item Montar y depurar el prototipo físico, analizando las variaciones respecto a la simulación.
\end{itemize}
\section{Marco Teórico}
\subsection{¿Qué es un Amplificador de Audio?}
Un amplificador de audio es un circuito electrónico cuya función principal es incrementar la amplitud de una señal de entrada variable en el tiempo (normalmente proveniente de fuentes como micrófonos, reproductores de audio o instrumentos) para entregar una potencia mayor a una cierta carga, tal como un altavoz o parlante.

\subsection{Clasificación de Amplificadores y Etapa de Potencia}

Los amplificadores se clasifican según su forma de amplificar la señal de entrada, separandose en clase A, B, AB, D, entre otros. El amplificador TDA2030 es de clase AB, el cual combina aspectos de los amplificadores de clase A y B.

Al igual que en el calse B, se utiliza un par de transistores push-pull, pero se elimina el curce por cero característico del amplificador de clase B. Este tipo de amplificador mejora el rendimiento típico de un clase A (al rededor del 25\%) a un 78\% aproximadamente, manteniendo la linealidad.

\section{Diseño en KiCad}
En esta sección se detalla el flujo de trabajo en el entorno EDA KiCad:
\subsection{Esquemático (\textit{Eeschema})}
Se ingresó el esquema eléctrico asignando los huellas (\textit{footprints}) correspondientes para cada componente pasivo y el encapsulado TO-220 del TDA2030.

\subsection{Diseño del PCB (\textit{Pcbnew})}
Para garantizar un trazado óptimo y reducir interferencias, se aplicaron los siguientes criterios:
\begin{itemize}[label=\small\color{primary}$\blacksquare$]
    \item Pistas de alimentación y salida hacia el parlante con ancho incrementado para soportar corrientes pico elevados.
    \item Planos de masa (GND) para minimizar el ruido inducido.
    \item Disposición física adecuada para permitir el montaje del disipador de calor térmico.
\end{itemize}
\section{Simulación Circuital}
Se realizaron simulaciones en entorno SPICE (o el simulador integrado en KiCad/LTspice) para evaluar:
\begin{enumerate}
    \item \textbf{Análisis de Respuesta en Frecuencia (AC Analysis):} Obtención del ancho de banda ($f_L$ y $f_H$) a $-3$\,dB.
    \item \textbf{Análisis Transitorio (Transient Analysis):} Verificación de la amplitud máxima de salida antes de alcanzar el recorte o saturación (\textit{clipping}).
\end{enumerate}
\section{Montaje Físico}
Se describe el proceso de ensamble sobre la placa de circuito impreso:
\begin{itemize}[label=\small\color{primary}$\blacksquare$]
    \item Inserción y soldadura de componentes (comenzando por componentes de menor perfil como resistencias y diodos, hasta los condensadores electrolíticos y el TDA2030).
    \item Montaje del disipador térmico con pasta de silicona conductora y aislante de mica/silicona.
    \item Inspección visual y prueba de continuidad previos a la energización.
\end{itemize}
\section{Depuración y Funcionamiento}
\subsection{Pruebas de Laboratorio}
Se conectó la placa a una fuente de alimentación regulada y se aplicó un generador de funciones en la entrada, midiendo la salida mediante un osciloscopio y una carga resistiva de prueba ($4\,\Omega$ / $8\,\Omega$).

\subsection{Resultados Obtenidos y Mediciones}
\begin{itemize}[label=\small\color{primary}$\blacksquare$]
    \item \textbf{Ganancia medida:} [Insertar valor medido $V_{out}/V_{in}$].
    \item \textbf{Potencia máxima sin distorsión:} $P_{out} = \frac{V_{rms}^2}{R_L} = \dots$
    \item \textbf{Comportamiento térmico:} Observaciones sobre la temperatura de trabajo con el disipador seleccionado.
\end{itemize}
\section{Conclusiones}
El amplificador basado en el TDA2030 demostró cumplir satisfactoriamente con las especificaciones del fabricante bajo el esquema estándar. Se constataron las bondades de la simulación previa para prever el comportamiento real, así como la importancia de un correcto trazado de PCB y un dimensionado térmico adecuado en circuitos de potencia de audio.
% ... tus inputs o secciones ...

\end{document}